\section*{2.5. The spectral theorem for compact symmetric operators}
\begin{proof}
We set $\alpha = \|A\|$ and assume $\alpha 
eq 0$ (i.e, $A 
eq 0$) without loss of
generality. Since
\[
\|A\|^2 = \sup_{\|f\|=1} \|Af\|^2 = \sup_{\|f\|=1} \langle Af, Af \rangle = \sup_{\|f\|=1} \langle f, A^2 f \rangle
\tag{2.44}
\]
there exists a normalized sequence $u_n$ such that
\[
\lim_{n\to\infty} \langle u_n, A^2 u_n \rangle = \alpha^2.
\tag{2.45}
\]
Since $A$ is compact, it is no restriction to assume that $A^2 u_n$ converges, say
$\lim_{n\to\infty} A^2 u_n = \alpha^2 u$. Now
\begin{align*}
\| (A^2 - \alpha^2)u_n \|^2 &= \| A^2 u_n\|^2 - 2\alpha^2 \langle u_n, A^2 u_n \rangle + \alpha^4 \\
&\leq 2\alpha^2(\alpha^2 - \langle u_n, A^2 u_n \rangle)
\tag{2.46}
\end{align*}
(where we have used $\| A^2 u_n\| \leq \|A\|\|A u_n\| \leq \|A\|^2\|u_n\| = \alpha^2$) implies
$\lim_{n\to\infty}(A^2 u_n - \alpha^2 u_n) = 0$ and hence $\lim_{n\to\infty} u_n = u$. In addition, $u$ is
a normalized eigenvector of $A^2$ since $\langle (A^2 - \alpha^2)u, u \rangle = 0$. Factorizing this last
equation according to $(A - \alpha)u = v$ and $(A + \alpha)v = 0$ show that either
$v 
eq 0$ is an eigenvector corresponding to $-\alpha$ or $v = 0$ and hence $u 
eq 0$ is an
eigenvector corresponding to $\alpha$.
\end{proof}
Note that for a bounded operator $A$, there cannot be an eigenvalue with
absolute value larger than $\|A\|$, that is, the set of eigenvalues is bounded by
$\|A\|$ (Problem 2.2).

Now consider a symmetric compact operator $A$ with eigenvalue $\alpha_1$ (as
above) and corresponding normalized eigenvector $u_1$. Setting
\begin{align*}
\tilde{\mathcal{H}}_1 &= \{u_1\}^\perp = \{f \in \tilde{H} \mid \langle u_1, f \rangle = 0\} \\
\mathcal{D}(A|_1) &= \mathcal{D}(A) \cap \tilde{\mathcal{H}}_1 = \{f \in \mathcal{D}(A) \mid \langle u_1, f \rangle = 0\}
\end{align*}
we can restrict $A$ to $\tilde{\mathcal{H}}_1$ using
\[
\langle u_1, Af \rangle = \langle Au_1, f \rangle = \alpha_1 \langle u_1, f \rangle = 0
\tag{2.47}
\]
since $f \in \mathcal{D}(A|_1)$ implies $\langle u_1, f \rangle = 0$ and hence $Af \in \tilde{\mathcal{H}}_1$. Denoting this restriction by $A|_1$, it is not hard to
see that $A|_1$ is again a symmetric compact operator. Hence we can apply
Theorem 2.17 iteratively to obtain a sequence of eigenvalues $\alpha_j$ with corresponding normalized eigenvectors $u_j$. Moreover, by construction, $u_n$ is
orthogonal to all $u_j$ with $j < n$ and hence the eigenvectors $\{u_j\}$ form an
orthonormal set. This procedure will not stop unless $\tilde{H}$ is finite dimensional.
However, note that $\alpha_j = 0$ for $j > n$ might happen if $A^n = 0$.

\begin{theorem}{2.18.}
Suppose $\tilde{H}$ is a Hilbert space and $A : \tilde{H} \to \tilde{H}$ is a compact
symmetric operator. Then there exists a sequence of real eigenvalues $\alpha_j$
\end{theorem}